% ===========================================================================
%  sviluppatore/04-motore/04-ensemble.tex
% ===========================================================================
\section{Comporre più stream}\label{sez:ensemble}

Una strategia da sola è uno stream. \simbolo{EnsembleEngine}, in
\file{reporting/ensemble.py}, ne prende diversi e ne fa uno solo. Il percorso
ha tre stadi, e due di essi sono punti d'innesto per i plugin.

\figuraTikz{composizione-ensemble}
  {La composizione. Pesatura e overlay fanno cose diverse e non vanno
   confuse: la prima distribuisce, la seconda scala.}
  {composizione-ensemble}

\subsection{Primo stadio: portare gli stream alla stessa scala}

Due strategie con lo stesso Sharpe ma volatilità diversa non vanno pesate
allo stesso modo: la più volatile dominerebbe il risultato
indipendentemente dal peso dichiarato. Ogni stream viene quindi riportato a
una volatilità obiettivo comune prima di essere pesato.

Il mattone è uno solo, e serve sia al backtest sia all'operatività dal vivo:

\begin{codiceInterno}{python}{Il mattone comune}{reporting/ensemble.py:335}
    def _vol_trailing(self, returns: pd.Series) -> pd.Series:
        """Vol annualizzata TRAILING e NON shiftata: il valore a t usa le
        osservazioni fino a t INCLUSO.

        E' il mattone comune fra backtest e live. Il backtest la shifta di un
        giorno (`_compute_scale_series`) perche' lo scale applicato al ritorno
        di t deve conoscere solo il passato; il live la usa non shiftata
        (`scala_corrente`) perche' li' l'ultima osservazione e' il ritorno
        appena realizzato e la posizione che si sta per aprire rendera' nel
        periodo SUCCESSIVO. Due usi, una sola formula: e' l'unico modo perche'
        il libro che si esegue sia quello che si e' misurato.
        """
\end{codiceInterno}

\begin{perche}{Perché uno shift di differenza non è un dettaglio}
Nel backtest il fattore applicato al rendimento di $t$ deve conoscere solo il
passato, quindi la serie si shifta di uno. Dal vivo l'ultima osservazione
\emph{è} il rendimento appena realizzato, e la posizione che si sta per
aprire renderà nel periodo successivo: shiftare anche lì significherebbe
dimensionare con una stima vecchia di un giorno. Stesso numero, due usi, e
se le due strade calcolassero la volatilità ciascuna per conto proprio il
libro eseguito non sarebbe quello misurato. Il motivo per cui la funzione è
una sola è esattamente questo.
\end{perche}

\noindent
Due costanti governano i casi limite:

\begin{center}
\begin{tabular}{@{}L{32mm}R{14mm}L{92mm}@{}}
\toprule
\textbf{Costante} & \textbf{Valore} & \textbf{A che serve} \\
\midrule
\rigaTenue
\cod{SCALE\_MIN\_OBS} & \num{20} & Sotto questa soglia non c'è stima, e il
  fattore vale \cod{1.0}: il libro resta quello dell'allocatore, mai
  amplificato da una volatilità stimata su quattro dati. \\
\cod{SCALE\_MAX\_LEVA} & \num{10.0} & Con volatilità che tende a zero il
  fattore esploderebbe. La volatilità ha quindi un pavimento a
  \cod{target\_vol / SCALE\_MAX\_LEVA}. \\
\bottomrule
\end{tabular}
\end{center}

\subsection{Secondo stadio: la pesatura}

Quattordici modalità installate, raggruppate per famiglia:

\begin{center}
\begin{tabular}{@{}L{34mm}L{18mm}L{86mm}@{}}
\toprule
\textbf{Famiglia} & \textbf{Tipo} & \textbf{Modalità} \\
\midrule
\rigaTenue
Ingenua & statica & \cod{equal} \\
Rischio / covarianza & statica & \cod{inverse\_cov}, \cod{min\_variance\_lw},
  \cod{min\_variance\_rmt}, \cod{erc} \\
\rigaTenue
Robusta & statica & \cod{hrp}, \cod{nco}, \cod{cluster\_risk\_budget},
  \cod{max\_div} \\
Coda & statica & \cod{cvar} \\
\rigaTenue
Con vista sui rendimenti & statica & \cod{black\_litterman}, \cod{kelly} \\
Adattiva al regime & \textbf{rolling} & \cod{regime\_ml}, \cod{regime\_style} \\
\bottomrule
\end{tabular}
\end{center}

\noindent
Il campo \cod{kind} del \cod{META} governa il dispatch, e non è
documentazione: è la condizione di un \cod{if}.

\begin{codiceInterno}{python}{Il dispatch, senza una catena di if}{reporting/ensemble.py:491}
        spec = next((m for m in self.WEIGHTING_MODE_REGISTRY
                     if m["key"] == self.weighting_mode), None)
        if spec is None:
            raise ValueError(f"Unknown weighting_mode: {self.weighting_mode}")
        weight_fn = _WEIGHTING_COMPUTERS[self.weighting_mode]
        if spec["kind"] == "static":
            weights = weight_fn(self, stream_pnls)
        else:  # rolling: pesi time-varying; per la UI esponi la MEDIA come scalare
            weights_over_time = weight_fn(self, stream_pnls)
            mean_w = weights_over_time.mean(axis=0)
            weights = {name: float(mean_w[name]) for name in stream_pnls}
\end{codiceInterno}

\begin{attenzione}{Il modello documenta un valore che non esiste}
\file{weights/WEIGHT\_MODE\_TEMPLATE.md} dice che \cod{kind} vale
\cod{"static"} o \cod{"dynamic"}. Nel codice il confronto è con
\cod{"static"}, e \emph{tutto il resto} finisce nel ramo rolling. I
quattordici plugin installati usano \cod{static} (dieci) e \cod{rolling}
(quattro); \cod{dynamic} non lo usa nessuno. Un plugin scritto seguendo il
modello funzionerebbe per caso. Il modello è da correggere, il codice no.
\end{attenzione}

\subsection{Terzo stadio: gli overlay}

Ogni overlay attivo produce un moltiplicatore in $(0,1]$ calcolato dal
PnL \emph{base}, cioè da quello prima di qualunque overlay. I moltiplicatori
si moltiplicano fra loro.

\begin{codiceInterno}{python}{L'applicazione degli overlay}{reporting/ensemble.py:555}
        # Overlay di esposizione CAUSALI (ortogonali al weighting_mode e componibili
        # fra loro). Ogni overlay attivo produce un moltiplicatore in (0,1] calcolato
        # dal PnL/stream BASE (=> ordine-indipendente); l'esposizione finale e' il
        # PRODOTTO dei moltiplicatori. Nessun overlay attivo => ensemble_pnl invariato.
        base_pnl = ensemble_pnl
        overlay_multipliers: dict[str, pd.Series] = {}
        for ov in self.OVERLAY_REGISTRY:
            if ov["key"] in self.enabled_overlays:
                m = _GATE_COMPUTERS[ov["key"]](self, base_pnl, stream_pnls)
                m = m.reindex(base_pnl.index).fillna(1.0)
                overlay_multipliers[ov["key"]] = m
        if overlay_multipliers:
            total_mult = pd.Series(1.0, index=base_pnl.index)
            for m in overlay_multipliers.values():
                total_mult = total_mult * m
            ensemble_pnl = base_pnl * total_mult
\end{codiceInterno}

\begin{nota}{L'indipendenza dall'ordine è una proprietà costruita}
Se ogni overlay calcolasse il proprio moltiplicatore dal PnL già modificato
dai precedenti, il risultato dipenderebbe dall'ordine in cui compaiono nel
registro --- cioè dal campo \cod{order}, che è un dato di presentazione.
Calcolandoli tutti dal PnL base, l'ordine smette di contare e attivare due
overlay dà lo stesso risultato comunque li si elenchi.
\end{nota}

\subsection{I sette overlay installati}

\begin{center}
\begin{tabular}{@{}L{26mm}L{32mm}L{80mm}@{}}
\toprule
\textbf{Chiave} & \textbf{Famiglia} & \textbf{Quando taglia} \\
\midrule
\rigaTenue
\cod{vol\_regime} & Volatilità & Quando la volatilità realizzata sale. Mira
  al drawdown massimo. \\
\cod{jump\_vol} & Volatilità & Reazione in un giorno a salti oltre $N$ sigma:
  copre il rischio di apertura che gli overlay lenti perdono. \\
\rigaTenue
\cod{corr\_gate} & Correlazione & Quando la correlazione media fra gli stream
  schizza, cioè quando la diversificazione sparisce. \\
\cod{regime\_gmm} & Stress & Miscela gaussiana su feature multivariate;
  riduzione graduale al crescere della probabilità di stress. \\
\rigaTenue
\cod{iso\_stress} & Stress & Isolation forest su dispersione, volatilità,
  ampiezza, correlazione media. \\
\cod{macro\_regime} & Macro & Fattore latente (PCA causale) su curva, spread
  di credito, VIX, disoccupazione. \\
\rigaTenue
\cod{news\_pressure} & Flusso informativo & Quando il volume di notizie
  supera la norma storica. È l'unico che dichiara dati esterni e una
  scadenza giornaliera. \\
\bottomrule
\end{tabular}
\end{center}
